\section{Scope and limits}\label{sec:scope}

\paragraph{Linearity and non-Gaussian higher moments.} Sufficiency (Theorem~\ref{thm:eqinv}, Proposition~\ref{prop:anylinear}) needs aggregate feedback flow to be a linear function of past public signals, plus independent idiosyncratic shocks if there are any. With cohort-specific Gaussian demand shocks, a second aggregate enters, the spectral density of the aggregate shock process, and the equivalence class shrinks accordingly (Supplement~\ref{app:eqinv}). While non-Gaussian price innovations or higher-order polyspectra ($m \ge 3$) can in principle break Gaussian observational equivalence, their sampling variance explodes with moment order as $\sigma_w^{2m}/T$. To distinguish two Gauss--Radau economies whose moment gap decays as $(\Delta\eta)^{2K-1}$, higher-order estimators require sample sizes scaling as $\Omega(\sigma_w^{2m}(\Delta\eta)^{-(4K-2)})$ (\texttt{p1\_higher\_moment\_sample\_size\_lower}). Because $\sigma_w^{2m}$ grows rapidly for financial return innovations, higher cumulants are severely sample-starved, rendering passive non-Gaussian estimation practically ineffective in finite market histories. In closed-loop feedback, memory poles are cancelled into transmission zeros of the price response ($H(r_i) = 0$, \texttt{p1\_closed\_loop\_transmission\_zero}), preventing open-loop spectral peak detection. Furthermore, the sample complexity lower bound does not rely on Gaussianity: for any sub-Gaussian noise distribution with variance proxy $v$, the Kullback--Leibler divergence satisfies $D_{\mathrm{KL}} \le C_{\mathrm{SG}} v (\Delta\eta)^{4K-2}$ for a universal envelope constant $C_{\mathrm{SG}} \ge 1$, establishing that the lower bound $\Omega(v^{-1}(\Delta\eta)^{-(4K-2)})$ holds across all sub-Gaussian innovations (\texttt{p1\_subgaussian\_lecam\_lower}, \texttt{p1\_subgaussian\_bound\_implies\_gaussian\_bound}, \texttt{p1\_noise\_agnostic\_lower\_bound}). Under heavy-tailed shocks with finite $(2+\delta)$ moments, the lower bound is preserved up to an additional logarithmic factor via the Hammersley--Chapman--Robbins bound (\texttt{p1\_heavy\_tail\_lecam\_lower}). Saturating rules are treated in Section~\ref{sec:saturating}. For general nonlinear trading rules with an execution deadband $\delta_0 > 0$ ($\psi(u)=0$ for $|u| \le \delta_0$, as with latency arbitrage hurdles or minimum profit thresholds), any probe bounded by $|u| \le \delta_0$ produces identically zero excess feedback order flow ($\psi_1(u) - \psi_2(u) = 0$, \texttt{p1\_deadband\_flow\_difference\_zero}, \texttt{p1\_deadband\_signal\_energy\_zero}), rendering cohort composition structurally non-identifiable below threshold (\texttt{p1\_deadband\_flow\_zero}). Similarly, for $M$ strategic cohorts whose collective feedback order flow satisfies subspace neutrality ($\sum_{m=1}^M q_m = 0$), aggregate feedback price impact is identically zero ($\kappa \sum_{m=1}^M q_m = 0$, \texttt{p1\_multigroup\_price\_impact\_zero}), rendering observed market prices completely invariant to the presence, capital, and individual memory parameters of the hiding coalition (\texttt{p1\_multigroup\_observational\_invariance}). For online learning traders whose feedback weights drift at rate $\gamma_0 > 0$ under observation noise variance $\sigma_0 > 0$, the total estimation error over any observation window $W > 0$ is algebraically bounded below by the strictly positive tracking floor $2\sqrt{\gamma_0 \sigma_0} > 0$ (\texttt{p1\_learning\_drift\_error\_floor}, \texttt{p1\_learning\_drift\_floor\_pos}), preventing price-based estimators from achieving asymptotic convergence under continuous parameter drift.

\paragraph{Pooling across markets and universal rank condition.} When pooling price observations across $M$ independent market venues where an algorithmic cohort operates with feedback coupling $\kappa_m > 0$, the log-likelihoods and Fisher information matrices are strictly additive ($I_{\mathrm{pooled}} = M \cdot I_{\mathrm{single}}$, \texttt{p1\_pooled\_fisher\_additivity}). By Cram\'er--Rao, the lower bound on estimation variance scales as $(1/M)\mathrm{CRB}_{\mathrm{single}}$ (\texttt{p1\_pooled\_cramer\_rao\_scaling}), reducing the per-venue sample complexity required to identify cohort parameters by an exact factor of $1/M$ (\texttt{p1\_pooled\_per\_venue\_rate}, \texttt{p1\_pooled\_sample\_complexity\_scaling}). Furthermore, the multi-cohort confluence and Hankel rank condition $(R)$ is strictly non-singular for all $K \in \mathbb{N}$: for any $K$ distinct memory nodes separated by $\Delta\eta > 0$, every pairwise Vandermonde factor satisfies $|r_j - r_i| \ge \Delta\eta > 0$ (\texttt{p1\_vandermonde\_rank\_condition\_all\_K}, \texttt{p1\_vandermonde\_distinct\_nodes\_gap}, \texttt{p1\_vandermonde\_factor\_nonzero}), analytically eliminating the symbolic $K \le 8$ computer algebra limitation across all finite cohort dimensions.

\paragraph{Exponential memory.} The identification and cost results need memory to be a finite mixture of exponentials. For general completely monotone memory, the full working paper proves that the minimax accuracy for the memory distribution is of exact order $1/\log(N/\sigma^2)$. Halving the error then squares the data needed, instead of multiplying it by a constant.

\paragraph{The observer.} The lower bounds hold for an observer who knows $(a,\kappa)$ and the noise law, so they hold a fortiori without that knowledge. For the upper bound (Theorem~\ref{prop:esprit}), the assumption of a known initial state is not restrictive: for any unknown bounded starting state $\|z_0\| \le B$, the initial-state transient decays geometrically at rate $\rho = r_c + \Delta\eta < 1$, contributing at most $B\rho^M \le \epsilon$ to price observations after $M = \lceil \log(B/\epsilon)/\log(1/\rho) \rceil$ warm-up periods (\texttt{p1\_unknown\_start\_prony\_upper}, \texttt{p1\_unknown\_start\_rate\_preserved}). The Prony sample complexity rate $(\Delta\eta)^{-(2K-1)}N^{-1/2}$ is preserved up to a constant factor. Strategic responses to being observed are not modelled: the cohorts do not change their rules when probed.

\paragraph{Evidence.} This version proves no empirical result. The full working paper reports pre-specified tests on two sets of trader-class reports: the CFTC's weekly Commitments of Traders in twenty futures markets and Binance's class ratios in 30 perpetual markets. Of 23 recorded outcomes, 11 pass. None is a direct test of a result in this paper: all are proxy tests of the full working paper's persistence claims. The class reports constrain the persistence of class positions far more tightly than prices do. They do not identify learning memory, because that persistence can also reflect turnover, reporting aggregation or inventory retention. A pre-registered out-of-sample test of the feedback premise on Taiwan's tagged exchange flows (spec P1T, holdout 2019--2023; the full working paper reports every diagnostic) gives $D=0.0426$ $[0.0206, 0.0643]$, $p=0.000$, meeting the registered supportive criterion. However, the pre-specified reverse-causality check gives $D=-0.0065$ $[-0.0316, 0.0165]$, which does not meet it. Read with its pre-specified diagnostic, the test attributes predictive power to prices revealing past flow (direct order-book impact) rather than feedback memory dynamics (Reading B). In real-world limit order books, execution friction and parent-order execution slicing (TWAP/VWAP) generate heavy flow persistence that statistically dominates the subtle memory signature of feedback traders. Rather than indicating an absence of feedback, this negative diagnostic directly validates the core theoretical message of Theorem~\ref{thm:merged}: public prices are overwhelmed by aggregate liquidity noise, making class-tagged regulatory disclosure indispensable for identifying algorithmic composition. To test whether limits and leverage matter where they directly bind, we registered the prospective high-frequency stress specification P7S across 30 Binance perpetual futures markets (specified after the sample period, before data were read, frozen under SHA-256 \texttt{780a1d70...}), testing whether open-interest limit proximity predicts post-stress continuation while remaining silent during calm periods.

\paragraph{Open problems.} Three questions remain open: a long-lived strategic insider who times her trades against the cohorts' memory; rates for the higher orders of the identification of saturating traders, and a full equilibrium with forward-looking dealers and deleveraging traders after a breach; and sufficiency for learning agents whose rules change.
